\providecommand{\main}{..} 
\documentclass[\main/boa.tex]{subfiles}

\begin{document}

\section{XGBoost rządzi}
\begin{minipage}[t]{0.915\textwidth}
	\center     
    \includegraphics[width=60px]{img/workshops/czarno_biale/vladimir-crop.png} 
\end{minipage}

\begin{minipage}{0.915\textwidth}
\centering
{\bf \index[a]{Alekseichenko Vladimir} Vladimir Alekseichenko}
\end{minipage}

\vskip 0.3cm

\begin{affiliations}
\begin{minipage}{0.915\textwidth}
\centering
\large General Electric \\[2pt]
\end{minipage}
\end{affiliations}

\vskip 0.8cm

\opiswarsztatu XGBoost to jest jedna najlepszych implementacji "Gradient Boosting" z punktu widzenia praktycznego.9
Dlaczego warto?
\begin{enumerate}
	\item Wynik (czyli zwykle jeden z najlepszych)
	\item Czas na naukę i predykcję (potrafi używać wszystkie dostępne rdzenie)
	\item Odporność na przeuczenia się (poprzez różne parametry regularyzacji)
	\item Stabilność (można spokojnie używać na produkcje)
\end{enumerate}

\planwarsztatu
\begin{enumerate}
\item Zrozumienie biznes problemu
\item Zrozumienie danych
\item Budowa bardzo prostego modelu (base-line)
\item Przypomnienie co to jest drzewa decezyjne
\item Uruchomienie prostego modelu xgboost
\item Generowanie cech (feature engineering)
\item Budowanie bardziej zaawansowanego modelu
\item Optymalizacja hyperparametrów
\item Inne (zaawansowane) triki (opcjonalnie)
\end{enumerate}	 

\end{document}
